\scanpage{43}

\margin{\figure{\begin{tikzpicture}[x=1cm,y=1cm]\draw (0.7,1.7)--(0,0)--(0.7,-1.7);\end{tikzpicture}}}

1) $V=V^{10}+V^{01}$ \unsure{décomposition de vectoriels} sur $C$

\deleted{1') $V^{10}\simeq\unsure{\check V^0}\otimes C,\quad V^{01}\simeq\unsure{\check V^1}\otimes_K C[-1]$}

2) Opération de $\pi$ sur $V$, \unsure{comp. aux op. de $\pi$ sur $C$}.

3) $V\simeq M\otimes_{\mathbb Z_p}C$, $M$ un \deleted{\unsure{vectoriel|réseau}} \inserted{module libre de t.f.} sur $\mathbb Z_p$.

\uline{Relations}

\uline{1) et 2)}
\[
V^{10}=V^\pi\otimes_K C
\]
\[
V^{01}=(V\otimes C[1])^\pi\otimes_K C[-1]
\]
\margin{\uline{NB} Structures $1,2$ \unsure{reviennent à un couple} $\unsure{\check V^0}=V^\pi$, $\unsure{\check V^1}=(V\otimes C(1))^\pi$.
NB. $2\unsure{\not\Rightarrow|\Rightarrow}1$}

\uline{2) et 3)} $M$ stable par $\pi$
\margin{[NB. \unsure{Structures $2,3$. connues} \inserted{\unsure{donne aussi 1)}} quand on \unsure{donne} $M$ avec op. de $\pi$ dessus
\deleted{\unsure{NB.} \illegible{2}} (i.e. le grp $p$-divisible \unsure{sur} $K$)]}

\uline{3 et 1} ??? \deleted{Les autom. de $M$ \unsure{qui sont compatibles}}

\uline{NB} le corps $K$ intervient dans 2), du moins \unsure{implicitement}, mais pas dans 1), ni 3), qui \unsure{semblent} des structures intrinsèques : la donnée du groupe $p$-divisible sur
\[
\unsure{\widehat S}\subset\widehat\Omega=C.
\]

Regardons le \inserted{plus petit} groupe \unsure{alg.} d'autom. \inserted{$G$} de $M\otimes Q_\ell$ \deleted{\unsure{qui}} tel que $G_C$ contienne les opérations
$\begin{pmatrix}\lambda&0\\0&\unsure{1|\xi}\end{pmatrix}$ pour la décomposition de $V$ en $V^{10}+V^{01}$.

\margin{53}
